\scanpage{87}
% archive foliation: 86
\diagram{\begin{tikzcd}[ampersand replacement=\&,row sep=large]
\text{Préschémas atomiques : (réunions de \unsure{spectres|schémas} de corps)}
\arrow[d,"\text{foncteur pleinement fidèle, commutant aux sommes}"']\\
\text{Foncteurs covariants }(\mathrm{Corps})\unsure{-|\to}(\mathrm{Ens})
\end{tikzcd}}
\margin{[Sous-catégorie pleine de $(\mathrm{Sch})$]}
i.e. contravariants $(\text{Schs ponctuels})^\circ\unsure{-|\to}(\mathrm{Ens})$

\diagram{\begin{tikzcd}[ampersand replacement=\&,row sep=large]
\text{Préschémas atomiques sur }k\ [\text{i.e. sur }\operatorname{Spec}(k)]
\arrow[d,"\text{foncteur pleinement fidèle}"']\\
\text{Foncteurs covariants }(\text{Extensions de }k)\unsure{-|\to}(\mathrm{Ens})
\end{tikzcd}}
\margin{$(\mathrm{At})/\operatorname{Spec}k$}
Foncteurs `` '' $(\text{Schs ponctuels}/\operatorname{Spec}k)^\circ\unsure{-|\to}(\mathrm{Ens})$

Soit $X$ un préschéma, il définit un préschéma atomique
\[
|X|=\coprod_{x\in X}\operatorname{Spec}k(x),
\]
qui \deleted{\unsure{est}} \inserted{\unsure{représente}}
la restriction du foncteur contravariant défini par
$X:(\mathrm{Sch})^\circ\unsure{-|\to}(\mathrm{Ens})$,
\deleted{\unsure{aux schémas ponctuels}}.
Bien sûr, on obtient un foncteur $X\mapsto|X|$,
dont la \uline{restriction aux préschémas réduits est fidèle}.
\deleted{\unsure{pleinement}}

Soient $X,Y$ deux préschémas, on appelle \uline{pseudo-morphisme}
ou \uline{application constructible} de $X$ dans $Y$, tt morphisme
$f:|\unsure{\widetilde X|\bar X}|\to|\unsure{\widetilde Y|\bar Y}|$
qui a la propriété suivante : Pour tte partie
\deleted{quasi-compact} \inserted{\unsure{ouvert affine}} $U$ de $X$,
il existe une partition de $U$ en des $U_i$ \inserted{loc. fermés}
(\uline{constructibles}), tels que $f||U_i|$ soit un
\uline{morphisme} $U_{i\,\mathrm{red}}\to Y$.

% Reversed blue typewriting is visible through the reused sheet.

% The reversed typewriting showing through this sheet is on the separately supplied reverse, scan 88, where it is transcribed.
